\section{Experiments}
\label{sec:experiment}


% Gray-cell helper for cited official results. Requires xcolor.
\newcommand{\g}[1]{\textcolor{gray}{#1}}

\begin{table*}[t]
\centering
\footnotesize
\setlength{\tabcolsep}{3pt}
\renewcommand{\arraystretch}{1.15}
\caption{\textbf{Results on WeaveBench.}
Gray rows denote official results reported by \cite{weavebench}, using the best reported thinking mode for each backbone; models other than GPT-5.5 and Claude Opus~4.7 were evaluated only with \textsc{OpenClaw}.
Black rows denote our runs with Qwen~3.7-Plus.
Our runs use root privileges inside the task virtual machine, whereas the official results use a regular user account and are therefore included as reference points rather than matched comparisons.
PR denotes full-task PassRate~(\%), and Overall denotes the mean score over all 114 tasks.
The remaining columns report PassRate across the eight benchmark domains.}

\label{tab:weavebench}
\begin{tabular}{l l cc | cccccccc}
\toprule
\textbf{Model} & \textbf{Harness} & \textbf{PR}$\uparrow$ & \textbf{Overall}$\uparrow$ & \textbf{DSK} & \textbf{DOC} & \textbf{GAM} & \textbf{WEB} & \textbf{DAV} & \textbf{OPS} & \textbf{SPA} & \textbf{DES} \\
\midrule
\multirow{4}{*}{\g{Claude Opus 4.7}}
& \g{\textsc{Claude Code}}  & \g{41.2} & \g{0.532} & \g{55.6} & \g{47.1} & \g{23.5} & \g{53.3} & \g{23.1} & \g{50.0} & \g{33.3} & \g{40.0} \\
& \g{\textsc{OpenClaw}}     & \g{35.1} & \g{0.482} & \g{55.6} & \g{29.4} & \g{23.5} & \g{66.7} & \g{15.4} & \g{41.7} & \g{16.7} & \g{20.0} \\
& \g{\textsc{Hermes Agent}} & \g{28.1} & \g{0.516} & \g{33.3} & \g{47.1} & \g{11.8} & \g{26.7} & \g{30.8} & \g{50.0} & \g{8.3}  & \g{10.0} \\
& \g{\textsc{Codex CLI}}    & \g{13.2} & \g{0.378} & \g{16.7} & \g{11.8} & \g{11.8} & \g{6.7}  & \g{7.7}  & \g{25.0} & \g{16.7} & \g{10.0} \\
\midrule
\multirow{4}{*}{\g{GPT-5.5}}
& \g{\textsc{Codex CLI}}    & \g{35.1} & \g{0.499} & \g{38.9} & \g{29.4} & \g{23.5} & \g{53.3} & \g{15.4} & \g{50.0} & \g{58.3} & \g{10.0} \\
& \g{\textsc{OpenClaw}}     & \g{33.3} & \g{0.466} & \g{38.9} & \g{35.3} & \g{35.3} & \g{21.4} & \g{23.1} & \g{38.5} & \g{33.3} & \g{40.0} \\
& \g{\textsc{Hermes Agent}} & \g{31.6} & \g{0.466} & \g{55.6} & \g{29.4} & \g{35.3} & \g{40.0} & \g{7.7}  & \g{25.0} & \g{25.0} & \g{20.0} \\
& \g{\textsc{Claude Code}}  & \g{14.9} & \g{0.299} & \g{33.3} & \g{11.8} & \g{11.8} & \g{0.0}  & \g{15.4} & \g{16.7} & \g{25.0} & \g{0.0}  \\
\midrule
\g{GPT-5.4}           & \g{\textsc{OpenClaw}} & \g{22.8} & \g{0.465} & \g{55.6} & \g{35.3} & \g{5.9}  & \g{0.0} & \g{23.1} & \g{23.1} & \g{8.3} & \g{20.0} \\
\g{GPT-5.3-codex}     & \g{\textsc{OpenClaw}} & \g{18.4} & \g{0.456} & \g{33.3} & \g{23.5} & \g{29.4} & \g{0.0} & \g{7.7}  & \g{16.7} & \g{8.3} & \g{20.0} \\
\g{GPT-5.2-codex}     & \g{\textsc{OpenClaw}} & \g{6.1}  & \g{0.321} & \g{5.6}  & \g{11.8} & \g{0.0}  & \g{0.0} & \g{15.4} & \g{16.7} & \g{0.0} & \g{0.0}  \\
\g{GPT-5.1-codex}     & \g{\textsc{OpenClaw}} & \g{1.8}  & \g{0.226} & \g{0.0}  & \g{5.9}  & \g{0.0}  & \g{0.0} & \g{7.7}  & \g{0.0}  & \g{0.0} & \g{0.0}  \\
\g{Gemini 3.1 pro}    & \g{\textsc{OpenClaw}} & \g{1.8}  & \g{0.223} & \g{0.0}  & \g{0.0}  & \g{0.0}  & \g{0.0} & \g{0.0}  & \g{8.3}  & \g{8.3} & \g{0.0}  \\
\g{Qwen3.5-397B-A17B} & \g{\textsc{OpenClaw}} & \g{0.9}  & \g{0.318} & \g{0.0}  & \g{0.0}  & \g{0.0}  & \g{0.0} & \g{0.0}  & \g{8.3}  & \g{0.0} & \g{0.0}  \\
\g{Qwen3-VL-8B-Think} & \g{\textsc{OpenClaw}} & \g{0.9}  & \g{0.092} & \g{0.0}  & \g{0.0}  & \g{0.0}  & \g{0.0} & \g{8.3}  & \g{0.0}  & \g{0.0} & \g{0.0}  \\
\g{GUI-Owl-1.5-32B}   & \g{\textsc{OpenClaw}} & \g{0.0}  & \g{0.065} & \g{0.0}  & \g{0.0}  & \g{0.0}  & \g{0.0} & \g{0.0}  & \g{0.0}  & \g{0.0} & \g{0.0}  \\
\midrule
\multirow{2}{*}{Qwen 3.7-Plus}
& Claude Code & 51.8 & 0.702 & 83.3 & 76.5 & 29.4 & 46.7 & 53.8 & 66.7 & 16.7 & 20.0 \\
& \textbf{LongHorizon-Harness (Claude Code)} & \textbf{80.7} & \textbf{0.835} & \textbf{88.9} & \textbf{100.0} & \textbf{58.8} & \textbf{73.3} & \textbf{84.6} & \textbf{91.7} & \textbf{66.7} & \textbf{80.0} \\
\bottomrule
\end{tabular}
\end{table*}

\subsection{Experimental Setup}
\label{sec:exp-setup}
\paragraph{Compared configurations.}
We use Qwen-3.7-Plus as the primary backbone model, and additionally evaluate Claude Opus-4.7 to assess the generality across backbone models. For each backbone model, we compare the original Claude Code evaluation framework~\cite{claudecode} with \textbf{LongHorizon-Harness}, where the latter uses Claude Code as its execution backend. Unless otherwise specified, the manager, executor, and auditor are instantiated with the same backbone model, ensuring that the comparison reflects the effect of task-state management rather than differences in model capability.

\paragraph{Implementation details.}
LongHorizon-Harness integrates existing agent backends through a unified AgentAdapter interface. Each role is assigned an independent execution budget: the executor is limited to 1800 seconds per round, while both the manager and the auditor are limited to 300 seconds. We set the maximum number of MEA rounds to $N_{\max}=25$. 
% The framework records the inputs and outputs of all roles, interaction trajectories, audit results, and runtime metadata, enabling subsequent analysis and reproducibility.

\paragraph{Benchmarks and metrics.}
We evaluate LongHorizon-Harness on three recent benchmarks that cover complementary forms of long-horizon execution.

\textbf{WeaveBench}~\cite{weavebench} consists of 114 tasks that require coordinated GUI and CLI interactions within the same workflow. We follow the standard WeaveBench evaluation protocol, using the task definitions, workspaces, runtime assets, and judge templates released by the official benchmark for all tasks. We report PassRate, defined as the percentage of fully passed tasks, and Overall, defined as the average score across all tasks. In addition, we report PassRate for each of the eight domains covered by the benchmark.


\textbf{OSWorld 2.0}~\cite{osworld2} consists of 108 desktop workflow tasks, with a median human completion time of approximately 1.6 hours. We use the official osworld-v2-2026.06.24 release and its standard Docker-based VM infrastructure. We report two metrics: Binary and Partial Accuracy. Binary Accuracy counts a task as successful only when its final score is 1, while Partial Accuracy is the average score over all tasks.


\textbf{Terminal-Bench 2.1}~\cite{terminalbench} is designed to evaluate challenging and realistic command-line tasks. We follow the evaluation setup used for the Qwen-series models, and evaluate on Terminal-Bench 2.1 using Harbor with the Docker backend. Final results report the average score over three independent runs per task, and are compared against the official leaderboard results.



\subsection{Main Results}
\label{sec:exp-results}

\paragraph{Effectiveness on long-horizon computer use.}
Table~\ref{tab:weavebench} reports the results on WeaveBench. The Claude Code harness with Qwen~3.7-Plus achieves a PassRate of 51.8\% and a mean task score of 0.702. Using the same model and retaining Claude Code as the executor backend, LongHorizon-Harness increases these results to 80.7\% and 0.835, respectively. This matched comparison isolates the contribution of the additional task-state management layer. Performance also improves across all eight task domains, showing that the gain is not concentrated in a particular application category. Since the officially reported configurations use a different privilege setting, we include them as reference points but base our main conclusion on the matched Claude Code comparison.

\begin{table}[t]
\centering
\footnotesize
\setlength{\tabcolsep}{4pt}
\renewcommand{\arraystretch}{1.12}
\caption{\textbf{Results on OSWorld~2.0.}
Gray rows denote official results reported by \cite{osworld2} under batched- and single-action settings, while the black row denotes LongHorizon-Harness with Qwen~3.7-Plus in the hybrid setting.
Binary is the percentage of tasks whose final benchmark score equals 1, and Partial is the mean benchmark score over
all 108 tasks. Bold marks the best official result.}
\label{tab:osworld}

\resizebox{0.8\linewidth}{!}{%
\begin{tabular}{l l cc}
\toprule
\textbf{Model} & \textbf{Harness / Mode} & \textbf{Binary}$\uparrow$ & \textbf{Partial}$\uparrow$ \\
\midrule
\g{Claude Opus 4.8} & \g{Batched actions} & \g{\textbf{20.6}} & \g{\textbf{54.8}} \\
\g{Claude Opus 4.7} & \g{Batched actions} & \g{18.2} & \g{48.9} \\
\g{GPT-5.5} & \g{Batched actions} & \g{13.0} & \g{49.5} \\
\midrule
\g{Claude Opus 4.8} & \g{Single action} & \g{18.5} & \g{49.3} \\
\g{Claude Opus 4.7} & \g{Single action} & \g{13.9} & \g{49.1} \\
\g{Claude Sonnet 4.6} & \g{Single action} & \g{8.3} & \g{41.5} \\
\g{MiniMax M3} & \g{Single action} & \g{4.6} & \g{22.3} \\
\g{Kimi 2.6} & \g{Single action} & \g{4.6} & \g{22.1} \\
\g{Qwen 3.7-Plus} & \g{Single action} & \g{2.8} & \g{21.5} \\
\midrule
Qwen 3.7-Plus & LongHorizon-Harness (hybrid) & 8.3 & 35.2 \\
\bottomrule
\end{tabular}%
}
\end{table}

\begin{figure}[t]
\centering
\includegraphics[width=0.9\linewidth]{pic/terminal_bench_rank_figure.png}
\caption{\textbf{Results on Terminal-Bench~2.1.}
LongHorizon-Harness retains Claude Code as its executor backend and improves success rate from 69.7\% to 77.2\%. Values marked with $^\ast$ are externally reported metrics.}
\label{fig:tb21}
\end{figure}

\begin{table}[t]
\centering
\footnotesize
\setlength{\tabcolsep}{5pt}
\renewcommand{\arraystretch}{1.15}
\caption{\textbf{OSWorld~2.0 Opus~4.7 subset (34 tasks).}
Both rows use Claude Opus~4.7 as the backbone.}
\label{tab:osworld-opus}

\resizebox{0.8\linewidth}{!}{%
\begin{tabular}{l l cc}
\toprule
\textbf{Model} & \textbf{Harness / Mode} & \textbf{Binary}$\uparrow$ & \textbf{Partial}$\uparrow$ \\
\midrule
Claude Opus 4.7 & \textsc{Single action}          & 20.6          & 55.8          \\
Claude Opus 4.7 & \textbf{LongHorizon-Harness} (hybrid) & \textbf{35.3} & \textbf{66.9} \\
\bottomrule
\end{tabular}%
}
\end{table}


Table~\ref{tab:osworld} reports the results on OSWorld~2.0. With Qwen~3.7-Plus, LongHorizon-Harness increases binary completion from 2.8\% to 8.3\% and partial score from 21.5\% to 35.2\%. The higher partial score shows that the agent satisfies a larger fraction of the task requirements, while the threefold increase in binary completion shows that this progress more frequently leads to a fully completed workflow. Together, the results on WeaveBench and OSWorld~2.0 show that explicit task-state management improves both cross-interface execution and sustained progress over long visual computer-use trajectories.


\paragraph{Complementarity with backbone capability.}
Table~\ref{tab:osworld-opus} further evaluates LongHorizon-Harness with Claude Opus~4.7 on a 34-task subset of OSWorld~2.0. It increases binary completion from 20.6\% to 35.3\% and partial score from 55.8\% to 66.9\%. The consistent improvements with Qwen~3.7-Plus and Claude Opus~4.7 show that the framework does not merely compensate for a weaker model. The backbone determines the quality of the actions and decisions produced within each round, while LongHorizon-Harness determines how reliably verified outcomes are maintained and used across rounds. Stronger models and explicit task-state management therefore provide complementary gains.


\paragraph{Generalization beyond computer use.}
Fig.~\ref{fig:tb21} reports the results on Terminal-Bench~2.1. Due to the limited budget, we only run three configurations ourselves: Claude Code with Qwen~3.7-Plus, LongHorizon-Harness with Claude Code and Qwen~3.7-Plus, and LongHorizon-Harness with Codex and GPT-5.6 Luna. All other results are quoted from the official Terminal-Bench~2.1 leaderboard or the corresponding model's official report.
LongHorizon-Harness improves the success rate of Qwen~3.7-Plus with Claude Code from 69.7\% to 77.2\%. With Codex as the executor backend, LongHorizon-Harness further reaches 83.1\% using GPT-5.6 Luna. Terminal-Bench tasks are performed entirely through the command line and require neither visual perception nor GUI--CLI routing. The gain therefore cannot be attributed to mechanisms specific to computer-use agents. It shows that explicit task-state maintenance, bounded execution contexts, and independent auditing also benefit general long-horizon agent execution. Across hybrid-interface, desktop, and pure CLI environments, LongHorizon-Harness consistently improves how agents preserve and convert intermediate progress into completed tasks.

\subsection{Analysis}
\label{sec:exp-analysis}

We next examine the computational cost of these gains, the task conditions under which they arise, and how the model and harness jointly determine the capability of the resulting agent system.

\paragraph{Cost--performance trade-off.}
Fig.~\ref{fig:osworld-frontier} shows that LongHorizon-Harness moves Qwen~3.7-Plus to a substantially stronger point on the OSWorld~2.0 cost--performance frontier. The same model improves from 2.8\% to 8.3\% in binary completion and from 21.5\% to 35.2\% in partial score, although its average output-token consumption increases from 28.9K to 104K per task. This result shows that a stronger harness can raise the task-level performance achievable by a fixed model. The increased cost observed with Qwen~3.7-Plus is specific to this model--task configuration; as analyzed below, the token cost of LongHorizon-Harness depends strongly on the capability of the underlying executor.

\begin{figure*}[t]
\centering
\includegraphics[width=1\linewidth]{pic/tokens_effort_figure.png}
\caption{\textbf{Cost--performance frontier on OSWorld~2.0.}
Binary completion (left) and partial score (right) versus average output tokens per task.
Colored curves show the official results under different reasoning-effort settings~\cite{osworld2}, with marker size indicating the reasoning effort.
The blue star denotes LongHorizon-Harness with Qwen~3.7-Plus, and the dashed arrow shows its improvement over the official single-action Qwen~3.7-Plus baseline.}
\label{fig:osworld-frontier}
\end{figure*}

\begin{figure*}[t]
\centering
\includegraphics[width=\linewidth]{pic/role_tokens.pdf}
\caption{\textbf{Token consumption by role.}
Average tokens per task consumed by the manager, executor, and auditor in LongHorizon-Harness, compared with the corresponding baseline.
Segment labels report the fraction consumed by each role.
All LH-Harness roles use Qwen~3.7-Plus.
WeaveBench and Terminal-Bench~2.1 report total tokens, whereas OSWorld~2.0 reports output tokens because the official provides only output-token statistics~\cite{osworld2}.}
\label{fig:role-tokens}
\end{figure*}



\paragraph{Computation across roles.}
Fig.~\ref{fig:role-tokens} decomposes the token consumption of LongHorizon-Harness across the manager, executor, and auditor. The manager accounts for only 2.8\%, 2.0\%, and 8.1\% of the total tokens on WeaveBench, OSWorld~2.0, and Terminal-Bench~2.1, respectively, showing that explicit task-state maintenance introduces little computational overhead. The auditor accounts for 19.4\%, 24.8\%, and 38.1\% of the tokens, indicating that independent verification constitutes the main additional investment of the framework. The overall token change nevertheless differs across benchmarks. LongHorizon-Harness consumes $2.3\times$ the baseline tokens on WeaveBench and $3.6\times$ the baseline output tokens on OSWorld~2.0, whereas it consumes 24\% fewer tokens on Terminal-Bench~2.1 while achieving a higher success rate. These results show that the framework does not impose a fixed token multiplier. Its total cost depends on both the task and the number of execution and recovery rounds required by the underlying model.


\begin{figure*}[ht]
\centering
\includegraphics[width=\linewidth]{pic/domain_deltas.pdf}
\caption{\textbf{Performance gains across task types.}
LongHorizon-Harness (blue) is compared with the same-model baseline (gray) across WeaveBench domains (left), OSWorld~2.0 capability tags (middle), and Terminal-Bench~2.1 categories (right).
The rightmost column of each panel reports the absolute performance change in points.
Blue connectors indicate improvements, and red connectors indicate regressions.
Parenthesized values report the number of OSWorld tasks or Terminal-Bench trajectories in each category.}
\label{fig:domain-deltas}
\end{figure*}

\begin{table*}[h]
\centering
\footnotesize
\setlength{\tabcolsep}{5pt}
\renewcommand{\arraystretch}{1.12}
\caption{\textbf{Per-task scores and token consumption on the WeaveBench Games subset} (17 tasks).
CC = Claude Code; LH = LongHorizon-Harness wrapping Claude Code as its backend.
Tok denotes total tokens per task (M).
\textbf{Bold} marks the better score per backbone.}
\label{tab:games}
\begin{tabular}{l cc cc | cc cc}
\toprule
& \multicolumn{4}{c}{\textbf{Claude Opus 4.7}} & \multicolumn{4}{c}{\textbf{Qwen 3.7-Plus}} \\
\cmidrule(lr){2-5} \cmidrule(lr){6-9}
& \multicolumn{2}{c}{Score} & \multicolumn{2}{c}{Tok (M)} & \multicolumn{2}{c}{Score} & \multicolumn{2}{c}{Tok (M)} \\
\cmidrule(lr){2-3} \cmidrule(lr){4-5} \cmidrule(lr){6-7} \cmidrule(lr){8-9}
\textbf{Task} & CC & LH & CC & LH & CC & LH & CC & LH \\
\midrule
\texttt{mines\_solve}           & \textbf{0.958} & 0.920 &  8.8 &  4.9 & 0.040 & \textbf{0.590} &  6.0 & 97.2 \\
\texttt{game\_ui\_bug}          & 0.830 & \textbf{0.910} & 11.7 &  8.3 & 0.750 & \textbf{0.850} & 10.0 & 34.2 \\
\texttt{stockfish\_puzzle}      & \textbf{0.603} & 0.550 &  7.5 &  3.4 & 0.000 & \textbf{0.550} &  0.9 & 15.2 \\
\texttt{gdb\_pygame\_cheat}     & 0.920 & \textbf{0.950} & 33.9 &  5.6 & 0.720 & \textbf{0.860} &  5.9 & 13.9 \\
\texttt{mines\_visual}          & 0.800 & \textbf{0.860} & 29.1 &  4.3 & 0.000 & \textbf{0.470} &  3.7 & 50.1 \\
\texttt{quadrapassel\_autoplay} & 0.600 & \textbf{0.870} & 50.5 & 13.4 & 0.000 & \textbf{0.300} & 10.0 & 59.1 \\
\texttt{pokerth\_equity\_play}  & 0.050 & \textbf{0.400} &  4.6 & 23.1 & 0.000 & \textbf{0.920} &  9.0 & 16.2 \\
\texttt{sokoban\_solver}        & \textbf{0.990} & 0.950 &  4.4 &  3.3 & \textbf{0.960} & 0.958 &  3.9 &  4.2 \\
\texttt{rhythm\_autoplay}       & 0.880 & \textbf{0.900} & 10.5 &  6.9 & \textbf{0.950} & 0.830 &  8.5 & 19.6 \\
\texttt{gnuchess\_blunder\_hunt}& 0.720 & \textbf{0.920} & 20.1 & 14.0 & 0.830 & \textbf{0.940} & 19.5 & 21.4 \\
\texttt{anagramarama\_word\_grid}& 0.550 & \textbf{0.750} &  6.6 & 18.3 & 0.750 & \textbf{0.830} & 12.8 & 12.5 \\
\texttt{godot\_scene\_debug}    & \textbf{0.958} & 0.920 & 11.4 &  7.8 & \textbf{0.930} & 0.870 &  8.6 & 17.1 \\
\texttt{hedgewars\_mission\_debug}& 0.790 & \textbf{0.810} & 19.6 &  8.5 & 0.740 & \textbf{0.830} & 14.4 & 13.1 \\
\texttt{supertux\_level\_repair}& 0.600 & \textbf{0.720} & 19.7 & 11.8 & 0.000 & \textbf{0.650} &  7.5 & 21.0 \\
\texttt{xmoto\_level\_repair}   & \textbf{0.600} & 0.550 &  8.1 & 23.4 & \textbf{0.920} & 0.850 & 19.4 & 83.3 \\
\texttt{kmahjongg\_pair\_solver}& 0.350 & \textbf{0.830} &  8.8 & 13.2 & \textbf{0.790} & 0.750 & 30.1 & 35.4 \\
\texttt{pysol\_freecell}        & 0.360 & \textbf{0.940} & 24.8 & 18.0 & \textbf{0.520} & 0.410 & 12.5 & 68.8 \\
\midrule
\textbf{Mean}                   & 0.680 & \textbf{0.809} & 16.5 & 11.1 & 0.524 & \textbf{0.733} & 10.7 & 34.3 \\
\bottomrule
\end{tabular}
\end{table*}


\paragraph{Task-dependent effectiveness.}
Fig.~\ref{fig:domain-deltas} compares the gains of LongHorizon-Harness across task types. On WeaveBench, Design and Spatial/3D obtain the largest PassRate improvements of 60.0 and 50.0 points, respectively, whereas the Desktop domain, which already reaches 83.3\% with the baseline, improves by 5.6 points. On OSWorld~2.0, some of the largest gains appear in streaming interaction, human-in-the-loop, and tutorial-following tasks. On Terminal-Bench~2.1, system-administration and game tasks improve substantially, while several short analytical categories show smaller gains or regressions. Despite operating through different interfaces, the categories with larger improvements commonly require agents to preserve, inspect, and revise multiple dependent environment states over an extended trajectory. The smaller improvements on several analytical categories suggest that LongHorizon-Harness provides less benefit when task performance is dominated by an individual model capability, such as visual perception, mathematical reasoning, coding, or algorithm design. Independent auditing can detect an incorrect result and initiate recovery, but it cannot supply a capability that the model does not possess. The effectiveness of LongHorizon-Harness is thus determined less by whether a task uses a GUI or CLI than by whether its primary bottleneck lies in long-horizon execution reliability or in solving an individual step.


\paragraph{Agent capability as a system property.}
Table~\ref{tab:games} compares the Claude Code baseline and LongHorizon-Harness on the same 17 WeaveBench Games tasks using Claude Opus~4.7 and Qwen~3.7-Plus. LongHorizon-Harness improves the mean score of Opus from 0.680 to 0.809 and that of Qwen from 0.524 to 0.733, with the largest gains concentrated on tasks where the corresponding baseline performs poorly. All six tasks on which the Qwen baseline scores at or below 0.04 recover to scores between 0.30 and 0.92, showing that the framework primarily raises the failure floor by recovering trajectories that would otherwise end in near-total failure. Qwen with LongHorizon-Harness further reaches a mean score of 0.733, exceeding the 0.680 obtained by Opus with the base Claude Code harness. This comparison demonstrates that agent capability is a property of the complete model--harness system. The model determines the actions and solutions available within each execution round, while the harness determines how reliably these local capabilities are decomposed, verified, recovered, and accumulated into end-to-end task completion. A stronger harness can consequently raise the task-level performance achievable with a fixed model, even though it does not add new primitive capabilities to that model. The token results reveal the same dependence on both components: LongHorizon-Harness increases Qwen's average consumption from 10.7M to 34.3M tokens per task, but reduces Opus consumption from 16.5M to 11.1M. This opposite pattern suggests that a stronger model can satisfy subtask contracts with fewer audit--replan rounds, whereas a weaker model must spend additional inference on repeated execution and recovery. Model capability and harness design jointly determine both the achievable performance and the computational cost of a long-horizon agent.


\subsection{Case Studies}
\label{sec:case-studies}

We further compare paired trajectories in which the Claude Code
baseline and LongHorizon-Harness execute the same task using the same
Qwen~3.7-Plus model. These cases examine what persists between rounds
under the Manage-Execute-Audit loop. Unresolved failures and completed
progress are carried outside the execution history, completion claims
are checked before entering the task state, and unmet dependencies
remain available when the manager defines subsequent subtasks. The
examples also illustrate how persistent task state connects executors
that operate in separate, fresh contexts.


\paragraph{Recovering from a stalled interaction.}
Fig.~\ref{fig:case-webrtc} compares how the two systems proceed after
the same GUI interaction fails. The baseline recognizes that the
Wireshark ``Decode As'' dialog is unresponsive, but this observation
remains embedded in its growing execution history. It consequently
returns to the same interaction for more than 400 steps without
collecting the remaining task evidence. LongHorizon-Harness instead
records the failed interaction and the unresolved evidence gaps in the
task state. Later executors focus directly on the missing
simulcast-layer charts and packet-level evidence, increasing the score
from 0.59 to 0.92. By externalizing the stalled interaction,
LongHorizon-Harness preserves the progress already verified and
provides the next executor with only the unresolved requirements.
Recovery therefore resumes from the latest audited task state rather
than from the failed trajectory that produced it.


\paragraph{Auditing apparent completion.}
Fig.~\ref{fig:case-heading} presents a case in which the executor stops
at a visually plausible result that does not satisfy the full task
specification. The baseline directly edits the document XML and
terminates after the headings appear correct. However, the task
requires the styles to be applied through the LibreOffice workflow,
and the result therefore receives a score of 0.00.
LongHorizon-Harness performs the prescribed GUI operations, confirms
the style changes during execution, and rechecks the document page by
page. The auditor then parses the document XML and confirms that all
15 headings have the required underlying style, yielding a score of
0.89. This independent audit prevents the executor's plausible but
non-compliant completion claim from becoming part of the persistent
task state. Subsequent planning can instead proceed from the document
state confirmed in the environment, rather than from the executor's
assessment of its own result.


\paragraph{Preserving pre-repair evidence.}
Fig.~\ref{fig:case-vlookup} presents a task in which part of the
required evidence belongs to the environment state before the repair.
The baseline captures an initial view of the formula errors, but
modifies the spreadsheet before completing the full pre-repair evidence
sequence. Its direct XML edit also leaves formula errors after the file
is reopened, producing an inconsistent before-and-after record.
LongHorizon-Harness records the missing pre-repair evidence as a
pending requirement. The next executor first completes this evidence,
then repairs the spreadsheet and checks the resulting formula cells.
The auditor subsequently reviews the complete sequence of nine
screenshots, increasing the score from 0.45 to 0.87. Keeping the
pre-repair evidence explicit changes the next subtask selected from the
task state: the original spreadsheet state is documented before the
repair is allowed to modify it. The resulting execution order preserves
a consistent evidence chain across the state transition.


\paragraph{Continuing from verified progress.}
Fig.~\ref{fig:case-lighthouse} compares two trajectories in which the
model can perform the central optimization but differs in how the
remaining workflow is maintained. The baseline improves the target
page, yet the same growing session must continue operating DevTools,
tracking the remaining deliverables, collecting evidence, and assessing
its own progress. Repeated DevTools interactions eventually exhaust
the available budget, leaving the task incomplete.
LongHorizon-Harness retains the completed optimization in the task
state and derives subsequent subtasks from the remaining requirements.
Fresh-context executors then confirm the page state, inspect the
relevant script structure, and collect the performance trace, while
the auditor checks the final screenshot against its metadata. The
score increases from 0.53 to 0.85. Because the completed optimization
persists outside the execution history, later executors can focus on
the remaining evidence without reconstructing or carrying the preceding
optimization process. The manager maintains continuity across the full
workflow, while each executor uses its context for the current
environment transition, allowing the model's local progress to be
carried forward into a more complete task result.


% Add this command to the preamble.
\DeclareRobustCommand{\taskid}[1]{\texttt{\detokenize{#1}}}


\begin{figure*}[h]
\centering
\includegraphics[width=\linewidth]{pic/case_webrtc_v3.pdf}
\caption[Recovering from a stalled interaction]{
\textbf{Recovering from a stalled interaction.}
The case is \taskid{WEB_task_16}, a WebRTC simulcast-layer audit.
\emph{Top:} The Claude Code baseline recognizes that Wireshark's
``Decode As'' dialog is unresponsive, but continues retrying the same
interaction for more than 400 steps and obtains a score of 0.59.
\emph{Bottom:} LongHorizon-Harness records the unresolved evidence gaps
in the task state. Subsequent execution rounds collect the missing
chart- and packet-level evidence, resulting in a score of 0.92.}
\label{fig:case-webrtc}
\end{figure*}

\begin{figure}[H]
\centering
\includegraphics[width=\linewidth]{pic/case_heading_v3.pdf}
\caption[Auditing apparent completion]{
\textbf{Auditing apparent completion.}
The case is \taskid{DOC_task_2}, which requires heading-style
normalization.
\emph{Top:} The Claude Code baseline edits the document XML directly
and terminates with a visually plausible result. Because the required
LibreOffice workflow is not followed, the result receives a score of
0.00.
\emph{Bottom:} LongHorizon-Harness applies the heading styles through
the prescribed GUI workflow and rechecks the document. The auditor then
parses the document XML to confirm the final style of all 15 headings,
obtaining a score of 0.89.}
\label{fig:case-heading}
\end{figure}


\begin{figure}[h]
\centering
\includegraphics[width=\linewidth]{pic/case_vlookup_v3.pdf}

\caption[Preserving pre-repair evidence]{
\textbf{Preserving pre-repair evidence.}
In \taskid{DOC_task_4}, the agent must repair a Calc VLOOKUP formula while preserving before-and-after evidence.
\emph{Top:} The Claude Code baseline records initial formula errors but edits the spreadsheet before completing the required pre-repair evidence, leaving the final file and evidence sequence inconsistent and scoring 0.45.
\emph{Bottom:} LongHorizon-Harness keeps the missing pre-repair evidence as a pending requirement, records it before modification, and audits the full nine-screenshot sequence, scoring 0.87.}
\label{fig:case-vlookup}
\end{figure}


\begin{figure}[H]
\centering
\includegraphics[width=\linewidth]{pic/case_lighthouse_v3.pdf}

\caption[Continuing from verified progress]{
\textbf{Continuing from verified progress.}
\taskid{WEB_task_10} is a Lighthouse performance-optimization task.
\emph{Top:} The Claude Code baseline completes the core optimization but gets stuck during DevTools evidence collection, fails to produce all required deliverables, and scores 0.53.
\emph{Bottom:} LongHorizon-Harness preserves the verified optimization state, delegates the remaining evidence requirements to later rounds, and audits the final screenshot and metadata, scoring 0.85.}
\label{fig:case-lighthouse}
\end{figure}